% SEGMENT OLP-0600-S01
\documentclass[../../../include/open-logic-section]{subfiles}

\begin{document}

\olfileid{sth}{choice}{vitali}
\olsection{பிற்சேர்க்கை: விட்டாலி முரண்புதிர்}

பனாக்--டார்ஸ்கி கட்டமைப்பு ஏற்கத்தக்கதா என்பதை நன்கு
உணர, அதன் \emph{நிறுவலை} ஆராய வேண்டும். ஆனால் அதற்கு
இங்கு வழங்க இயலாத அளவு இயற்கணிதம் தேவை. இருப்பினும்
பின்வரும் முடிவை மையமாகக் கொண்டு, பனாக்--டார்ஸ்கி
நிறுவலைப் புரிந்துகொள்ள உதவும் சில சுருக்கமான கருத்துகளைக்
கூறலாம்.\footnote{விரிவான விளக்கத்திற்கு \cite{Weston2003}
அல்லது \cite{Wagon2016} ஐப் பார்க்கவும்.}

\begin{thm}[விட்டாலி முரண்புதிர் ($\ZFC$ இல்)]\ollabel{vitaliparadox}
எந்த வட்டவரையையும் எண்ணத்தக்க எண்ணிக்கையிலான துண்டுகளாகப்
பிரித்து, அவற்றைச் சுழற்றியும் இடம் மாற்றியும் மீண்டும்
அமைத்தால், அந்த வட்டவரையின் இரு பிரதிகளை உருவாக்கலாம்.
\end{thm}

விட்டாலி முரண்புதிரை நிறுவுவது பனாக்--டார்ஸ்கி முரண்புதிரை
நிறுவுவதைவிட எளிது. அளவிட முடியாத கணம் பற்றிய
\citeyear{Vitali1905} இல் விட்டாலி வழங்கிய கட்டமைப்பிலிருந்து
இது கிடைப்பதால், இதை ``விட்டாலி முரண்புதிர்'' என்கிறோம்.
ஆனால் இரு நிறுவல்களின் கணக் கோட்பாட்டுப் பகுதி மிகவும்
ஒத்தது. முக்கிய வேறுபாடு: பனாக்--டார்ஸ்கி \emph{முடிவுறு}
துண்டாக்கலைக் கருதுகிறது; விட்டாலி முரண்புதிர்
\emph{எண்ணத்தக்க முடிவிலி} துண்டாக்கலைக் கருதுகிறது.
\citet{Weston2003} கூறுவது போல, விட்டாலி முரண்புதிர்
பனாக்--டார்ஸ்கி அளவுக்கு வியப்பூட்டாதது; ஆனால் ``எளிய''
நிலைகளில்கூட வடிவியல் முரண்புதிர்கள் நிகழ முடியும் என்பதை
அது காட்டுகிறது.

% SEGMENT OLP-0600-S02
விட்டாலி முரண்புதிர் இருபரிமாண உருவமான வட்டவரையைப்
பற்றியது. ஆகவே $\Real^2$ என்ற தளத்தில் வேலை செய்வோம்.
% SOURCE CORRECTION TA-STH-047: சுழற்சிக் கோணங்கள் வெறும் பகுத்தெண் ரேடியன்கள் அல்ல; π இன் பகுத்தெண் மடங்குகள்.
$\rotationsgroup$ என்பது ஆதிப்புள்ளியைச் சுற்றி கடிகாரத்
திசையில், $[0,2\pi)$ இடைவெளியிலுள்ள பையின்
\emph{பகுத்தெண் மடங்குகளான} ரேடியன் கோணங்களால் புள்ளிகளைச்
சுழற்றும் சார்புகளின் கணம் என்க. $\rotationsgroup$ பற்றிய
சில இயற்கணித உண்மைகள் இவை (கூற்றின் பொருள் உடனே
புரியாவிட்டால், நிறுவல் அதைத் தெளிவாக்கும்):

\begin{lem}\ollabel{rotationsgroupabelian}
சார்புக் கோப்பின்கீழ் $\rotationsgroup$ ஒரு பரிமாற்றுக்
{குழு} ஆகும்.
\end{lem}

\begin{proof}
$0$ ரேடியன் சுழற்சியை $0_{\rotationsgroup}$ என
எழுதுவோம். இது $\rotationsgroup$ இன் அடையாள உறுப்பு;
ஏனெனில் எந்த $\rho \in \rotationsgroup$ க்கும்
$\comp{0_{\rotationsgroup}}{\rho} = \comp{\rho}{0_{\rotationsgroup}} =
\rho$.

ஒவ்வொரு உறுப்புக்கும் எதிருறுப்பு உள்ளது. $\rho \in
\rotationsgroup$ என்பது $r$ ரேடியன்களால் சுழற்றினால்,
$\rho^{-1} \in \rotationsgroup$ என்பது $2\pi - r$
ரேடியன்களால் சுழற்றுகிறது; எனவே
$\rho \circ \rho^{-1} = 0_\rotationsgroup$.
அந்தக் கோணம் பூச்சியம் எனில் அதன் எதிருறுப்பு அடையாளச் சுழற்சியே.
பையின் பகுத்தெண் மடங்குகள் முழுச் சுற்று அளவில்
கூட்டலின்கீழும் மூடப்பட்டவை என்பதால் கோப்பும்
குழுவுக்குள்ளேயே உள்ளது.

கோப்பு சேர்ப்புப் பண்புடையது: எந்த
$\rho, \sigma, \tau \in \rotationsgroup$ க்கும்
$\comp{\rho}{(\comp{\sigma}{\tau})} =
\comp{(\comp{\rho}{\sigma})}{\tau}$.

கோப்பு பரிமாற்றுப் பண்புடையது: எந்த
$\rho, \sigma \in \rotationsgroup$ க்கும்
$\comp{\rho}{\sigma} = \comp{\sigma}{\rho}$.
\end{proof}

இந்தக் குழு $\rotationsgroup$ ஐ இரண்டு பாதிகளாகப்
பிரித்து, எந்தப் பாதியிலிருந்தும் முழுக் குழுவை மீண்டும்
பெறலாம்:

\begin{lem}\ollabel{disjointgroup}
$\rotationsgroup$ ஐ ஒன்றையொன்று வெட்டாத
$\rotationsgroup_{1}$, $\rotationsgroup_{2}$ என்ற இரு
கணங்களாகப் பிரிக்கலாம்; ஒவ்வொன்றும்
$\rotationsgroup$ முழுவதையும் தரும் அடித்தளமாகும்.
\end{lem}

\begin{proof}
$[0, \pi)$ இல் உள்ள பையின் பகுத்தெண் மடங்குக்
கோணங்களின் சுழற்சிகள் $\rotationsgroup_{1}$ இல் இருக்கட்டும்;
$\rotationsgroup_{2} = \rotationsgroup
\setminus \rotationsgroup_1$ என்க. தொடக்க இயற்கணிதத்தால்
$\Setabs{\comp{\rho}{\rho}}{\rho \in \rotationsgroup_1} =
\rotationsgroup$. இதேபோன்ற முடிவு
$\rotationsgroup_2$ க்கும் கிடைக்கும்.
\end{proof}

% SEGMENT OLP-0600-S03
குழுக்களைப் பற்றிய இந்த உண்மையால்
\olref{vitaliparadox} ஐ நிறுவுவோம்.
$\onesphere$ என்பது அலகு வட்டவரை; அதாவது தளத்தின்
ஆதிப்புள்ளியிலிருந்து சரியாக~$1$ அலகு தொலைவிலுள்ள
புள்ளிகளின் கணம்:
$\Setabs{\tuple{r,s} \in \Real^2}{\sqrt{r^2+s^2}=1}$.
$\onesphere$ மீது பின்வரும் தொடர்பைக் கருதி அதைத்
துண்டுகளாகப் பிரிப்போம்:
\[
	r \sim s \emph{ அப்போதும் அப்போதுதான் }(\exists \rho \in \rotationsgroup)\rho(r) = s.
\]
அதாவது, ஆதிப்புள்ளியைச் சுற்றியுள்ள பையின் பகுத்தெண்
மடங்குக் கோணச் சுழற்சியால் ஒரு புள்ளியிலிருந்து மற்றொன்றை
அடைய முடிந்தால், $\onesphere$ இன் அவ்விரு புள்ளிகளும்
இந்தத் தொடர்பில் உள்ளன. எதிர்பார்த்தபடி:

\begin{lem}
$\sim$ ஒரு சமவலுத் தொடர்பு.
\end{lem}

\begin{proof}
\olref{rotationsgroupabelian} இன்படி தெளிவு.
\end{proof}

இப்போது தேர்வைப் பயன்படுத்தி, $\sim$ கீழ்
$\onesphere$ இன் ஒவ்வொரு சமவலு வகுப்பிலிருந்தும் சரியாக
ஒரு உறுப்பைக் கொண்ட $C$ என்ற கணத்தைப் பெறுகிறோம்.
அதாவது சமவலு வகுப்புகளின் பின்வரும் கணத்தின் மீது
ஒரு தேர்வுச் சார்பு $f$ ஐக் கருதுகிறோம்.\footnote{
$\rotationsgroup$ !!{enumerable} என்பதால் $E$ இன்
ஒவ்வொரு !!{element} உம் !!{enumerable}. ஆனால்
$\onesphere$ !!{nonenumerable}; ஆகவே
\olref{vitalicover} மற்றும்
\olref[card-arithmetic][simp]{kappaunionkappasize} இன்படி
$E$ !!{nonenumerable}. எனவே இங்கு \emph{எண்ணத்தக்கதல்லாத}
தேர்வு பயன்படுகிறது.}
\[
	E = \Setabs{\equivrep{r}{\sim}}{r \in \onesphere},
\]
$C = \ran{f}$ என்க. ஒவ்வொரு சுழற்சி
$\rho \in \rotationsgroup$ க்கும்,
$\funimage{\rho}{C}$ என்பது $C$ இன் ஒவ்வொரு புள்ளியையும்
$\rho$ மூலம் சுழற்றிப் பெறும் புள்ளிகளின் கணம்.
இந்தக் கணங்கள் வட்டவரை முழுவதையும் இடைவெளியின்றியும்
மேற்படிதலின்றியும் மூடுவதை அடுத்த இரண்டு முடிவுகள்
காட்டுகின்றன:

\begin{lem}\ollabel{vitalicover}
$\onesphere = \bigcup_{\rho \in \rotationsgroup} \funimage{\rho}{C}$.
\end{lem}

\begin{proof}
$s \in \onesphere$ என்க. அப்போது
$r \in \equivrep{s}{\sim}$ ஆகும் ஏதேனும் $r \in C$
உள்ளது. அதாவது $r \sim s$; எனவே ஏதேனும்
$\rho \in \rotationsgroup$ க்கு $\rho(r) = s$.
\end{proof}

\begin{lem}\ollabel{vitalinooverlap}
$\rho_1 \neq \rho_2$ எனில்
$\funimage{\rho_1}{C} \cap \funimage{\rho_2}{C} = \emptyset$.
\end{lem}

\begin{proof}
$s \in \funimage{\rho_{1}}{C} \cap \funimage{\rho_{2}}{C}$
என்க. அப்போது ஏதேனும் $r_{1}, r_{2} \in C$ க்கு
$s = \rho_{1}(r_{1}) = \rho_{2}(r_{2})$.
எனவே $\rho^{-1}_2(\rho_1(r_1)) = r_2$, மேலும்
$\comp{\rho_1}{\rho^{-1}_2} \in \rotationsgroup$;
ஆகவே $r_{1} \sim r_{2}$. ஒவ்வொரு $\sim$-சமவலு
வகுப்பிலிருந்தும் $C$ சரியாக ஒரு உறுப்பையே தேர்வதால்
$r_1 = r_2$. எனவே $s = \rho_1(r_1) = \rho_2(r_1)$;
இதனால் $\rho_1 = \rho_2$.
\end{proof}

% SEGMENT OLP-0600-S04
முன்னைய இயற்கணித உண்மைகளை வட்டவரைக்குப் பயன்படுத்துவோம்:

\begin{lem}\ollabel{pseudobanachtarski}
$\onesphere$ ஐ ஒன்றையொன்று வெட்டாத $D_{1}$,
$D_{2}$ என்ற இரு கணங்களாகப் பிரிக்கலாம். $D_{1}$ ஐ
எண்ணத்தக்க எண்ணிக்கையிலான கணங்களாகப் பிரித்து, அவற்றைச்
சுழற்றி $\onesphere$ இன் ஒரு பிரதியை அமைக்கலாம்;
$D_{2}$ க்கும் இதேபோல் செய்யலாம்.
\end{lem}

\begin{proof}
\olref{disjointgroup} இன் $\rotationsgroup_{1}$,
$\rotationsgroup_{2}$ ஆகியவற்றைப் பயன்படுத்தி:
\begin{align*}
	D_{1} &= \bigcup_{\rho \in \rotationsgroup_1} \funimage{\rho}{C} &
	D_{2} &= \bigcup_{\rho \in \rotationsgroup_2} \funimage{\rho}{C}
\end{align*}
\olref{vitalicover} இன்படி இது $\onesphere$ இன்
ஒரு பிரிப்பு; \olref{vitalinooverlap} இன்படி $D_1$,
$D_2$ ஒன்றையொன்று வெட்டாதவை. கட்டமைப்பின்படி $D_1$
ஐ ஒவ்வொரு சுழற்சிக்கும் ஒரு கணமான
% SOURCE CORRECTION TA-STH-048: வரையறுக்கப்படாத R_1 க்கு பதில் முன் வரையறுத்த rotationsgroup_1 ஐப் பயன்படுத்துக.
$\funimage{\rho}{C}$, $\rho \in \rotationsgroup_1$ என்ற
எண்ணத்தக்க கணங்களாகப் பிரிக்கலாம். இவற்றைச் சுழற்றி
$\onesphere$ இன் பிரதியை அமைக்கலாம். ஏனெனில்
$\onesphere = \bigcup_{\rho \in
\rotationsgroup}\funimage{\rho}{C} = \bigcup_{\rho \in
\rotationsgroup_1}\funimage{(\comp{\rho}{\rho})}{C}$;
இது \olref{disjointgroup}, \olref{vitalicover}
ஆகியவற்றிலிருந்து கிடைக்கிறது. $D_2$ க்கும் இதே
காரணவாதம் பொருந்தும்.
\end{proof}\noindent
இதிலிருந்து விட்டாலி முரண்புதிர் உடனே கிடைக்கிறது.
$\onesphere$ ஐ எண்ணத்தக்க எண்ணிக்கையிலான
$\funimage{\rho}{C}$ துண்டுகளாகப் பிரித்து,
அவற்றை உருமாறா இயக்கங்களால் மாற்றுவதன் மூலம்
$\onesphere$ இலிருந்து $\onesphere$ இன் \emph{இரு} பிரதிகளைப்
பெறலாம். $D_1$ இன் ஒவ்வொரு துண்டையும் சுழற்றவும்;
$D_2$ இன் ஒவ்வொரு துண்டையும் முதலில் இடம் மாற்றி,
பிறகு சுழற்றவும்.

நிறுவலின் உத்தியை மீண்டும் கூறுவோம். தளச் சுழற்சிகளின்
குழுவைப் பற்றிய சில இயற்கணித உண்மைகளிலிருந்து
தொடங்கினோம். அதைக் கொண்டு $\onesphere$ ஐ சமவலு
வகுப்புகளாகப் பிரித்தோம். பின்னர் ஒவ்வொரு வகுப்பிலிருந்தும்
உறுப்பைத் தேர்ந்தெடுக்கத் தேர்வைப் பயன்படுத்தியதால்
``முரண்புதிர்'' கிடைத்தது.

% SEGMENT OLP-0600-S05
பனாக்--டார்ஸ்கியை நிறுவவும் இதே உத்தியைத்தான்
பயன்படுத்துகிறோம். முக்கிய வேறுபாடு: அதன் இயற்கணித
உண்மைகள் விட்டாலி முரண்புதிருக்கானவற்றைவிட மிகவும்
சிக்கலானவை. ஆனால் அந்த இயற்கணித உண்மைகளுக்குத்
தேர்வுடன் தொடர்பில்லை. அவற்றைச் சுருக்கமாகக் கூறுவோம்.

பனாக்--டார்ஸ்கிக்காக முதலில் \olref{disjointgroup}
போன்ற ஒரு முடிவை நிறுவுகிறோம்: எந்த
\emph{கட்டற்ற குழு} வையும் நான்கு பகுதிகளாகப் பிரித்து,
அவற்றை ``நகர்த்தி'' முழுக் குழுவின் இரு பிரதிகளைப்
பெறலாம்.\footnote{\emph{நான்கு} பகுதிகள் போதும்
என்பது \cite{Robinson1947} இன் முடிவு. அண்மைய
நிறுவலுக்கு \citet[Theorem 5.2]{Wagon2016} ஐப் பார்க்கவும்.
குழுவின் பகுதிகளை ``நகர்த்துதல்'' என்று விவரிப்பதில்
\citet[p.~3]{Weston2003} ஐப் பின்பற்றுகிறோம்.}
பின்னர் $\Real^3$ இன் ஆதிப்புள்ளியைச் சுற்றியுள்ள
இரண்டு குறிப்பிட்ட சுழற்சிகளால், சுழற்சிகளின்
$F$ என்ற கட்டற்ற குழுவை உருவாக்கலாம் என்று
காட்டுகிறோம்.\footnote{\citet[Theorem 2.1]{Wagon2016}
ஐப் பார்க்கவும்.} (இதுவரை தேர்வு பயன்படுத்தப்படவில்லை.)
இப்போது கோளப் பரப்பின் இரண்டு புள்ளிகளில் ஒன்று
$F$ இன் ஒரு சுழற்சியால் மற்றொன்றிலிருந்து பெறப்பட்டால்
அவை ``ஒத்தவை'' என்கிறோம். அவ்வாறு ``ஒத்த''
புள்ளிகளின் ஒவ்வொரு சமவலு வகுப்பிலிருந்தும் சரியாக
ஒரு புள்ளியைத் தேர்ந்தெடுக்க \emph{தேர்வைப் பயன்படுத்துகிறோம்}.
\olref{pseudobanachtarski} இல் செய்ததுபோல,
$F$ இன் பிரிப்பைக் கோளப் பரப்புக்குப் பயன்படுத்தி,
அதை நான்கு பகுதிகளாகப் பிரிக்கிறோம். அவற்றை
``நகர்த்தி'' கோளப் பரப்பின் இரு பிரதிகளைப் பெறலாம்.
இதுவே \citep{Hausdorff1914} ஐ நிறுவுகிறது:

\begin{thm}[ஹவுஸ்டார்ஃப் முரண்புதிர் ($\ZFC$ இல்)]
எந்தக் கோளத்தின் பரப்பையும் முடிவுறு எண்ணிக்கையிலான
துண்டுகளாகப் பிரித்து, அவற்றைச் சுழற்றியும் இடம்
மாற்றியும் மீண்டும் அமைத்தால், அக்கோளப் பரப்பின்
ஒன்றையொன்று வெட்டாத இரு பிரதிகளை உருவாக்கலாம்.
\end{thm}

முழுப் பனாக்--டார்ஸ்கி தேற்றத்தைப் பெற இன்னும் சில
இயற்கணித உத்திகள் தேவை; அது கோளப் பரப்பை மட்டுமல்ல,
அதன் உள்ளகத்தையும் பற்றியது. ஆனால் முக்கியமான
இயற்கணிதப் பணிக்கு மேலே சேரும் இறுதி அலங்காரமே அது.
இதனால்தான் வெஸ்டன் எழுதுகிறார்:
\begin{quote}
  [\ldots] கட்டற்ற குழுக்களைப் பற்றிய முடிவே
  பனாக்--டார்ஸ்கி முரண்புதிரின் நிறுவலில்
  \emph{முக்கியப் படி}. இந்தப் பார்வையில்,
  பனாக்--டார்ஸ்கி முரண்புதிர் $\Real^3$ பற்றிய
  கூற்றாக இருப்பதைவிட [$\Real^3$ இன் இடப்பெயர்ச்சிகளும்
  சுழற்சிகளும் அமைக்கும்] குழுவின் சிக்கல்தன்மை பற்றிய
  கூற்றாக இருக்கிறது. \cite[p.~16]{Weston2003}
\end{quote}
அதாவது \emph{முடிவுறு} துண்டாக்கல்
(பனாக்--டார்ஸ்கிபோல்) கிடைக்கிறதா, \emph{எண்ணத்தக்க
முடிவிலி} துண்டாக்கல்தான் (விட்டாலிபோல்) கிடைக்கிறதா
என்பது, இருபரிமாணத்திலும் முப்பரிமாணத்திலும் வேலை
செய்யும் குழுக்களின் சில இயற்கணித உண்மைகளைப் பொறுத்தது.

இது \olref[banach]{sec} இன் இறுதியில் சொன்ன
நகைச்சுவையை ஓரளவு கெடுக்கிறது. ``பனாக்--டார்ஸ்கி''
என்பது இருபரிமாண வடிவம் என்பதால், அதன் துண்டுகளை
மறுஅமைத்துப் ``பனாக்--டார்ஸ்கி பனாக்--டார்ஸ்கி''
என்று பெற வேண்டுமானால், எண்ணத்தக்க
\emph{முடிவிலி} துண்டுகள் தேவைப்படும். நகைச்சுவையைச்
சரிசெய்ய, முப்பரிமாணத்தில் எழுத வேண்டும். இதை
வாசகருக்குப் பயிற்சியாக விடுகிறோம்.

% SEGMENT OLP-0600-S06
இறுதியாக ஒரு கருத்து. \olref[banach]{sec} இல்,
கோளத்தின் ``துண்டுகள்'' \emph{அளவிடத்தக்கவை}
அல்ல; கற்பனைப் படமாக வரைய முடியாத ``முடிவின்றிச்
சிதறல்கள்'' என்று குறிப்பிட்டோம்.
\olref{pseudobanachtarski} ஐப் பெறத் தேர்வைப்
பயன்படுத்தியதிலும் இதுவே உண்மை. இதை விளக்குவது
பயனுள்ளது.

மீண்டும் சில பின்னணியைச் சுருக்கமாகக் கூறுவோம்
(இது \emph{வெறும்} சுருக்கம் மட்டுமே; \emph{அளவு}
பற்றிய பாடநூல் விளக்கத்தைப் படிக்கலாம்).
கணம் $X$ க்கு ஓர் அளவை வரையறுப்பது என்றால்,
$X$ மீதான ஓர் ``$\sigma$-இயற்கணிதம்'' இல் உள்ள
ஒவ்வொரு $E$ க்கும் $\mu(E) \in \Real$ என்ற மதிப்பை
ஒதுக்குவதாகும். அளவு எதிர்மறையற்றது. இங்குள்ள
விவரங்களைவிட முக்கியமானது $\mu$ இன்
எண்ணத்தக்க கூட்டற்பண்பு: ஒன்றையொன்று வெட்டாத கணங்களின்
எண்ணத்தக்க ஒன்றிப்பின் அளவு, அவற்றின் தனித்தனி
அளவுகளின் கூட்டுத்தொகை. அதாவது $X_n$ கள்
ஒன்றையொன்று வெட்டாதவை எனில்
$\mu(\bigcup_{n < \omega} X_n) = \sum_{n < \omega}\mu(X_n)$.
ஒரு கணம் ``அளவிட முடியாதது'' என்றால், இங்கு கருதும்
மாறாமைப் பண்புடைய அளவின்கீழ் அதற்குத் தேவையான
வகையில் மதிப்பை ஒதுக்க முடியாது என்று
பொருள். முன்புள்ள $\rotationsgroup$ ஐப் பயன்படுத்தி:

\begin{cor}[விட்டாலி]
% SOURCE CORRECTION TA-STH-050: அளவிடத்தக்க கணங்களின் களமும் சுழற்சிகளின்கீழ் மூடப்பட்டிருக்க வேண்டும்.
$\mu(\onesphere) = 1$ ஆகவும், $X$, $Y$ ஒத்துருவானவை
எனில் $\mu(X) = \mu(Y)$ ஆகவும் இருக்கும் ஓர் அளவு
$\mu$ என்க. அதன் அளவிடத்தக்க கணங்களின் களம்
இக்குழுச் சுழற்சிகளின்கீழ் மூடப்பட்டிருக்கட்டும்.
அப்போது எல்லா $\rho \in
\rotationsgroup$ க்கும் $\funimage{\rho}{C}$ அளவிட
முடியாதது.
\end{cor}

\begin{proof}
மறுதலித்து ஏதேனும் ஒரு துண்டு அளவிடத்தக்கது என்க.
அப்போது ஏதேனும் $\sigma \in \rotationsgroup$
மற்றும் $r \in \Real$ க்கு
$\mu(\funimage{\sigma}{C}) = r$.
% SOURCE CORRECTION TA-STH-049: ρ சுழற்சிக் குழுவில் ஓடுகிறது; தேர்வுக் கணம் C இல் அல்ல.
ஒவ்வொரு $\rho \in \rotationsgroup$ க்கும்
$\funimage{\rho}{C}$ மற்றும்
$\funimage{\sigma}{C}$ ஒத்துருவானவை;
ஆகவே ஒவ்வொரு $\rho \in \rotationsgroup$ க்கும்
$\mu(\funimage{\rho}{C}) = r$.
\olref{vitalicover}, \olref{vitalinooverlap} இன்படி
$\onesphere = \bigcup_{\rho \in \rotationsgroup}\funimage{\rho}{C}$
என்பது ஒன்றையொன்று வெட்டாத கணங்களின் எண்ணத்தக்க
ஒன்றிப்பு. எனவே எண்ணத்தக்க கூட்டற்பண்பின்படி,
$\mu(\onesphere) = 1$ என்பது எல்லா
$\funimage{\rho}{C}$ களின் அளவுகளின் கூட்டுத்தொகை:
\[
	1 = \mu(\onesphere) = \sum_{\rho \in \rotationsgroup}\mu(\funimage{\rho}{C}) = \sum_{\rho \in \rotationsgroup}r
\]
ஆனால் $r = 0$ எனில் $\sum_{\rho \in \rotationsgroup}r = 0$;
$r > 0$ எனில் $\sum_{\rho \in \rotationsgroup}r = \infty$.
இரண்டும் ஒன்று என்ற முன்னெண்ணத்துக்கு முரண்.
\end{proof}

\end{document}
